% ECONOMICS_PROBLEM: {"id": "C73-38", "title": "Classification and bifurcation of genuinely periodic MFG equilibria", "jel_code": "C73", "source_id": "OP-0217"}
% FIRST_POSTED: 2026-10-09
% ATTEMPT_STATUS: unresolved
% ENTRY_KIND: research_attempt
% !TEX program = pdflatex
\documentclass[11pt]{article}
\usepackage[T1]{fontenc}
\usepackage[utf8]{inputenc}
\usepackage{lmodern}
\usepackage[margin=1in]{geometry}
\usepackage{amsmath,amssymb,amsthm,mathtools}
\usepackage[colorlinks=true,linkcolor=blue,citecolor=blue,urlcolor=blue]{hyperref}
\allowdisplaybreaks
\newtheorem{theorem}{Theorem}
\newtheorem{proposition}[theorem]{Proposition}
\newtheorem{lemma}[theorem]{Lemma}
\newtheorem{corollary}[theorem]{Corollary}
\newcommand{\R}{\mathbb R}
\newcommand{\E}{\mathbb E}
\newcommand{\Prb}{\mathbb P}
\newcommand{\Q}{\mathbb Q}
\newcommand{\T}{\mathbb T}
\newcommand{\F}{\mathcal F}
\newcommand{\Ind}{\mathbf 1}
\DeclareMathOperator{\sgn}{sgn}
\DeclareMathOperator{\diver}{div}
\DeclareMathOperator{\Ent}{Ent}
\title{Periodic mean-field games:\\nonlinear exclusion and resonance for arbitrary bifurcation sequences}
\author{GPT 6 Astra Ultra --- Version 2}
\date{9 October 2026}
\begin{document}
\maketitle
\noindent\textbf{Problem ID:} \texttt{C73-38}. \textbf{Research attempt; independent review pending.}

\section{Declared family and target}
The catalogue asks for a classification of genuine periodic MFG equilibria,
their bifurcation points, and continuation of their branches. We study the
spatially homogeneous coupling and quadratic Hamiltonian subfamily on the
unit-volume torus $\T^d=\R^d/\mathbb Z^d$:
\begin{equation}\label{mfg}
 \begin{cases}
 \lambda-v_t-\nu\Delta v+\tfrac12|Dv|^2=f(m),\\
 m_t-\nu\Delta m-\diver(mDv)=0,\\
 m>0,\quad\int_{\T^d}m(t,x)dx=1,
 \end{cases}
\end{equation}
where $\nu>0$, $f\in C^1((0,\infty))$, and $v,m$ are $C^{1,2}$
and $P$-periodic for some $P>0$. Fix $\int_0^P\int v=0$.
The stationary branch is $(v,m,\lambda)=(0,1,f(1))$.
Genuine periodicity means that $(Dv,m)$ is nonconstant in time.

We retain the nonlinear, period-independent exclusion region and complete
linear mode calculation from Version 1. Version 2 proves a stronger
necessary resonance theorem for arbitrary finite-period bifurcation
sequences, with no differentiable-branch or nonzero-tangent assumption. All statements are proved below. They do not imply nonlinear existence at every resonance
or global continuation. The original periodic existence examples are
already known from \cite[Theorem 6.1 and Remark 6.4]{CN}; the published
agenda is discussed in \cite[Section 3.3]{CD}. Version 1 records a check of those sources and their scope on
9 October 2026; the present version retains those references and does not
claim an additional literature search. The computations here are not
claimed as new bifurcation theory.

\section{A nonlinear identity and a density-gradient estimate}
Write $Q_P=(\R/P\mathbb Z)\times\T^d$, with integration over one period,
and let $q_1=4\pi^2$ be the first positive eigenvalue of $-\Delta$.

\begin{lemma}\label{identities}
Every solution of \eqref{mfg} satisfies
\begin{align}
 0&=\int_{Q_P}(f(m)-f(1))(m-1)
       +\tfrac12\int_{Q_P}(m+1)|Dv|^2,\label{monoid}\\
 \nu\int_{Q_P}\frac{|Dm|^2}{m}
   &=-\int_{Q_P}Dm\cdot Dv.\label{entropyid}
\end{align}
In particular,
\begin{equation}\label{cauchy}
 \int_{Q_P}m|Dv|^2\geq\nu^2\int_{Q_P}\frac{|Dm|^2}{m}.
\end{equation}
\end{lemma}
\begin{proof}
Subtract the first stationary equation from the first equation of
\eqref{mfg}, multiply by $\rho=m-1$, and integrate over $Q_P$.
The constant multiplier $\lambda-f(1)$ disappears because $\int\rho=0$
at every time. Periodicity and the Fokker--Planck equation give
\[
 \int_{Q_P}(-v_t)\rho
 =\int_{Q_P}v\rho_t
 =-\nu\int_{Q_P}Dv\cdot D\rho-\int_{Q_P}m|Dv|^2.
\]
The diffusion term $-\nu\Delta v$ cancels the first term on the right.
Combining the remaining term with $|Dv|^2\rho/2$ proves
\eqref{monoid}.

For \eqref{entropyid}, multiply the Fokker--Planck equation by
$1+\log m$ and integrate. Positivity and compactness make this a valid
smooth multiplier. Its time derivative integrates to zero over a period;
spatial integration by parts gives exactly \eqref{entropyid}.
Let $I=\int|Dm|^2/m$ and $J=\int m|Dv|^2$.
Weighted Cauchy--Schwarz gives $\nu I\leq\sqrt{IJ}$.
If $I=0$, \eqref{cauchy} is automatic; otherwise division by
$\sqrt I$ and squaring prove it.
\end{proof}

\begin{theorem}[A nonlinear exclusion region]\label{exclude}
Let $0<\varepsilon<1$, $L\geq0$, and suppose
\[
 (f(s)-f(1))(s-1)\geq-L(s-1)^2
       \quad(1-\varepsilon\leq s\leq1+\varepsilon).
\]
If
\begin{equation}\label{threshold}
 L<\nu^2q_1\frac{2+\varepsilon}{2(1+\varepsilon)^2},
\end{equation}
then, for every $P>0$, any solution of \eqref{mfg} satisfying
$\|m-1\|_\infty\leq\varepsilon$ is the stationary solution
$(v,m,\lambda)=(0,1,f(1))$.
Consequently, if $f'(1)>-\nu^2q_1$, there is a neighborhood of $m=1$
in the uniform norm containing no nonstationary periodic solutions,
uniformly over all periods.
\end{theorem}
\begin{proof}
The ratio $(m+1)/(2m)$ is at least
$c_\varepsilon=(2+\varepsilon)/(2(1+\varepsilon))$.
Lemma~\ref{identities}, the bound $m\leq1+\varepsilon$, and spatial
Poincar\'e inequality for the mean-zero function $m-1$ give
\begin{align*}
 \tfrac12\int(m+1)|Dv|^2
 &\geq c_\varepsilon\int m|Dv|^2
 \geq c_\varepsilon\nu^2\int\frac{|Dm|^2}{m}\\
 &\geq\frac{c_\varepsilon\nu^2}{1+\varepsilon}\int|Dm|^2
 \geq\nu^2q_1\frac{2+\varepsilon}{2(1+\varepsilon)^2}
                              \int|m-1|^2.
\end{align*}
All integrals here are over $Q_P$. Substitution in \eqref{monoid} yields
\[
 0\geq\left(\nu^2q_1\frac{2+\varepsilon}{2(1+\varepsilon)^2}-L\right)
                              \int|m-1|^2.
\]
The coefficient is positive, so $m=1$. Identity \eqref{monoid} then
forces $Dv=0$. The HJB equation says $v_t=\lambda-f(1)$;
periodicity makes this constant derivative zero, and normalization gives
$v=0$.

For the final assertion choose
$\max\{0,-f'(1)\}<L<\nu^2q_1$.
Continuity of $f'$ ensures the secant inequality in a sufficiently small
density neighborhood. The factor on the right of \eqref{threshold}
tends to one as $\varepsilon\downarrow0$, so shrinking the neighborhood
also makes \eqref{threshold} hold. No step depends on $P$.
\end{proof}

\begin{corollary}[The homogeneous monotone region]
If $f$ is nondecreasing on $(0,\infty)$, then \eqref{mfg} has no
genuinely periodic solution of any amplitude or period; its only
normalized periodic solution is $(0,1,f(1))$.
\end{corollary}
\begin{proof}
Both integrands in \eqref{monoid} are nonnegative, so $Dv=0$.
The Fokker--Planck equation reduces to $m_t=\nu\Delta m$.
Multiplication by $m$ and integration over one period gives
$\int|Dm|^2=0$. Unit mass then implies $m=1$, and the HJB equation
and normalization determine $v$ and $\lambda$ as in the preceding proof.
\end{proof}

\section{All linear periodic modes and what they imply}
\begin{proposition}\label{spectrum}
Let $a=f'(1)$ and linearize \eqref{mfg} about $(0,1,f(1))$.
For a real spatial Laplace eigenfunction $\phi$ with
$-\Delta\phi=q\phi$, $q>0$, write
$w(t,x)=U(t)\phi(x)$ and $\rho(t,x)=R(t)\phi(x)$.
The mode equations are
\begin{equation}\label{matrix}
 \frac{d}{dt}\begin{pmatrix}U\\R\end{pmatrix}
 =A_q\begin{pmatrix}U\\R\end{pmatrix},\qquad
 A_q=\begin{pmatrix}\nu q&-a\\-q&-\nu q\end{pmatrix},
 \qquad A_q^2=q(\nu^2q+a)I.
\end{equation}
A nonconstant $P$-periodic mode exists if and only if
\begin{equation}\label{resonance}
 a=-\nu^2q-\frac{(2\pi n/P)^2}{q}
 \quad\hbox{for some integer }n\geq1.
\end{equation}
If $a=-\nu^2q$, the periodic modes at that spatial eigenvalue are
stationary and equal to the kernel of $A_q$. If $a> -\nu^2q$, there
is no nonzero periodic mode at that eigenvalue. The spatially constant
mode vanishes after mass and value normalization, including the variation
of $\lambda$.
\end{proposition}
\begin{proof}
The linear equations at positive spatial frequencies are
$-w_t-\nu\Delta w=a\rho$ and
$\rho_t-\nu\Delta\rho-\Delta w=0$. Substitution yields
\eqref{matrix}; direct multiplication gives its square.
If $q(\nu^2q+a)>0$, the eigenvalues are nonzero real numbers of opposite
sign, so neither exponential over a positive period equals one.
If that product is zero, $A_q\ne0$ and $A_q^2=0$, so
$e^{PA_q}=I+PA_q$. Periodic initial vectors are exactly $\ker A_q$
and give stationary modes.
If the product is negative, put
$\omega_q=\sqrt{-q(\nu^2q+a)}>0$. Then
\[
 e^{tA_q}=I\cos(\omega_qt)+\frac{A_q}{\omega_q}\sin(\omega_qt).
\]
Its eigenvalues over one period are $e^{\pm i\omega_qP}$.
A nonzero real fixed vector exists precisely when
$\omega_qP=2\pi n$ for a positive integer $n$, proving
\eqref{resonance}. At resonance every initial vector in this two-dimensional
mode is periodic. The spatial zero mode has $\rho_0=0$ by mass,
$-w_0'+\delta\lambda=0$ by the HJB equation, and hence
$\delta\lambda=0$ by periodicity; the normalization makes $w_0=0$.
Expansion in a complete real Fourier eigenbasis gives the stated list
of modes without assuming simple eigenvalues.
\end{proof}

\begin{proposition}[A necessary condition for a differentiable periodic branch]
Let $f_\eta$ be smooth in $(\eta,m)$, and suppose a branch of
periodic solutions that is $C^1$ as a curve in the rescaled $C^{1,2}$
function space, with parameters $(\eta_\epsilon,P_\epsilon)$,
passes through $(0,1,f_{\eta_0}(1))$ at $\epsilon=0$.
Take derivatives after scaling time to a common unit period, subtract the
stationary branch in the multiplier, and assume the resulting first
variation $(w,\rho)$ is nonconstant in time. Then some positive spatial
eigenvalue $q$ and integer $n\geq1$ satisfy \eqref{resonance} with
$a=f'_{\eta_0}(1)$ and $P=P_0$.
\end{proposition}
\begin{proof}
Differentiating the rescaled equations is legitimate by the assumed
$C^1$ convergence in $C^{1,2}$. Terms involving $P'_0$ multiply time
derivatives of the stationary state and vanish. The derivative of
$f_\eta(1)$ cancels against the subtracted stationary multiplier.
The remaining equations are exactly the linear equations of
Proposition~\ref{spectrum} with period $P_0$ after returning to physical
time. A time-dependent first variation has a time-dependent Fourier
mode, to which that proposition applies.
\end{proof}

For example, $q=q_1=4\pi^2$ gives the basic oscillatory frequency
$\omega_1=2\pi\sqrt{-4\pi^2\nu^2-a}$ when $a<-4\pi^2\nu^2$.
At $\nu=1$, its fundamental period
$1/\sqrt{-4\pi^2-a}$ matches the normalization in
\cite[Theorem 6.1]{CN}. That theorem supplies nonlinear existence in
its specified regime; a calculation of imaginary eigenvalues by itself
does not.


\section{A resonance test without a differentiable branch}
The last proposition of Version 1 required a differentiable branch with
nonstationary first variation. A branch may have a zero first variation,
or a stationary tangent, and a bifurcating sequence need not lie on any
chosen differentiable curve. We now remove these restrictions for the
homogeneous quadratic family at a finite limiting period. The proof
uses time differences of the actual nonlinear solution, so stationary
linear degeneracies do not enter its conclusion.

Fix $\eta_0$ and $P_0>0$. Let $f_\eta(s)$ be $C^2$ jointly on a
neighborhood of $(\eta_0,1)$ and write $a_\eta=\partial_s f_\eta(1)$.
Spatial Fourier frequencies are $\xi=2\pi k$, $k\in\mathbb Z^d$;
write $q=|\xi|^2$ and $\omega_n=2\pi n/P$. All $L^2$ norms below
use normalized measure $P^{-1}dt\,dx$ on $Q_P$.
The zero temporal mean condition for a function $r$ means
$\int_0^P r(t,x)dt=0$ for every $x$, or its $L^2$ interpretation.
The corresponding condition for a spatial gradient is imposed on that
gradient, not on an arbitrary spatially constant part of its primitive.

\begin{lemma}[A uniform multiplier away from oscillatory resonance]
\label{multiplier}
Suppose
\begin{equation}\label{nonresonance}
 a_{\eta_0}\ne-\nu^2q-\frac{(2\pi n/P_0)^2}{q}
 \quad\text{for every }q=4\pi^2|k|^2>0
                    \text{ and every integer }n\geq1.
\end{equation}
There are neighborhoods $J$ of $\eta_0$ and $I$ of $P_0$ and a constant
$K<\infty$ such that, for $\eta\in J$, $P\in I$, every periodic solution
of
\begin{equation}\label{forcedlinear}
 -w_t-\nu\Delta w-a_\eta r=h+\ell(t),\qquad
 r_t-\nu\Delta r-\Delta w=\diver J_0
\end{equation}
with $\int_{\T^d}r(t)=0$ and zero temporal means for $Dw,r$ obeys
\begin{equation}\label{multbound}
 \bigl(\|Dw\|_2^2+\|r\|_2^2\bigr)^{1/2}
 \leq K\bigl(\|h\|_2^2+\|J_0\|_2^2\bigr)^{1/2}.
\end{equation}
Here $J_0$ is a vector field and $\ell(t)$ is any spatially constant
function. The notation $J_0$ in this lemma is unrelated to the parameter
neighborhood $J$. Square-integrable right-hand sides suffice when the
equations hold distributionally and the left-hand quantities are defined.
No assumption excludes stationary degeneracies $a_{\eta_0}=-\nu^2q$.
\end{lemma}
\begin{proof}
The zero spatial frequency contributes nothing to $Dw$ or $r$.
The zero temporal frequencies of these quantities also vanish by
assumption. At every remaining frequency, write $a=a_\eta$,
$\omega=\omega_n$ and let hats denote Fourier coefficients.
The system is
\[
 \begin{pmatrix}\nu q-i\omega&-a\\q&\nu q+i\omega\end{pmatrix}
 \binom{\widehat w}{\widehat r}
   =\binom{\widehat h}{i\xi\cdot\widehat J_0}.
\]
Its determinant and inverse formulas are
\begin{align*}
 D_{q,n}(a,P)&=\nu^2q^2+\omega_n^2+a q,\\
 \widehat w&=\frac{(\nu q+i\omega)\widehat h
                         +a\,i\xi\cdot\widehat J_0}{D_{q,n}},\\
 \widehat r&=\frac{-q\widehat h
                     +(\nu q-i\omega)i\xi\cdot\widehat J_0}{D_{q,n}}.
\end{align*}
Condition \eqref{nonresonance} makes these determinants nonzero at
$(a_{\eta_0},P_0)$ for all $q>0$ and $n\ne0$.
We justify uniform boundedness over the entire frequency set, including
perturbations of the period. First restrict to compact initial
neighborhoods where $|a|\leq A$ and $0<P_-\leq P\leq P_+<\infty$.
If $q\geq2A/\nu^2$, then
\[
 D_{q,n}\geq\tfrac12\nu^2q^2+\omega_n^2.
\]
There are only finitely many spatial frequencies below that threshold.
For any such $q$, when $|n|$ is sufficiently large uniformly in
$P\in[P_-,P_+]$, one has $\omega_n^2\geq2Aq$, and hence
$D_{q,n}\geq\nu^2q^2+\tfrac12\omega_n^2$.
Thus outside a finite set of pairs of frequencies,
\[
 D_{q,n}\geq c_0(q^2+\omega_n^2)>0
\]
for a constant $c_0$ depending only on $\nu$ and the initial
neighborhoods. On the finite remaining set, nonresonance and continuity
allow the neighborhoods to be reduced so that all $|D_{q,n}|$ are
bounded away from zero.

Using $|\xi\cdot\widehat J_0|\leq\sqrt q|\widehat J_0|$,
the four scalar bounds for the inverse mapping into
$(\sqrt q|\widehat w|,|\widehat r|)$ have coefficients
\[
 \frac{\sqrt q\sqrt{\nu^2q^2+\omega^2}}{|D_{q,n}|},\quad
 \frac{|a|q}{|D_{q,n}|},\quad
 \frac{q}{|D_{q,n}|},\quad
 \frac{\sqrt q\sqrt{\nu^2q^2+\omega^2}}{|D_{q,n}|}.
\]
Each is uniformly bounded on the finite set. On its complement the
determinant estimate proves the same fact: for example the first is at
most a fixed constant times
$\sqrt q/(q^2+\omega^2)^{1/2}\leq C/\sqrt q$, and
$q\geq4\pi^2$. The other two coefficients are bounded by constants
times $q/(q^2+\omega^2)\leq1/q$.
Choose $K$ larger than the resulting uniform matrix norm.
Apply the coefficient bound, sum its squares over all frequencies, and
use Parseval's identity. This proves \eqref{multbound}. Any frequencies
of $h,J_0$ outside the relevant set only increase its right side. The
function $\ell(t)$ has zero coefficient at every positive spatial
frequency and therefore has no effect on this argument.
\end{proof}

\begin{theorem}[No small genuine periodic solutions off resonance]
\label{nonlinearresonance}
Under \eqref{nonresonance}, there are neighborhoods $J,I$ as in
Lemma~\ref{multiplier} and $\varepsilon_0>0$ such that, for every
$\eta\in J$ and $P\in I$, every $C^{1,2}$ periodic solution of
\eqref{mfg} with $f=f_\eta$ and
\[
 \|m-1\|_\infty+\|Dv\|_\infty\leq\varepsilon_0
\]
is stationary in time. It may be a nonuniform stationary solution; the
theorem does not exclude a stationary spatial bifurcation.
Consequently every genuine periodic bifurcation point $(\eta_0,P_0)$
from the homogeneous branch, in the catalogue's finite-period sense,
satisfies \eqref{resonance} for some $q>0$ and $n\geq1$.
This conclusion applies to arbitrary converging sequences and needs
neither a differentiable branch nor a condition on its first variation.
\end{theorem}
\begin{proof}
Put $r=m-1$, $g=Dv$, and $a=a_\eta$. Define the nonlinear remainders
\[
 H_\eta(r,g)=f_\eta(1+r)-f_\eta(1)-a r-\tfrac12|g|^2,
 \qquad J(r,g)=r g.
\]
The equations can be written
\[
 -v_t-\nu\Delta v-a r=H_\eta(r,g)+f_\eta(1)-\lambda,
 \qquad r_t-\nu\Delta r-\Delta v=\diver J(r,g).
\]
For a temporal shift $s$ let $\delta_s r(t)=r(t+s)-r(t)$, with
times taken modulo $P$, and similarly define $\delta_s v,g,H,J$.
The difference fields solve \eqref{forcedlinear} with
$w=\delta_s v$, density $\delta_s r$, and forcings $\delta_s H$,
$\delta_s J$; the scalar constants cancel. They have zero temporal
mean because integration over a full period is invariant under the shift.
Also $\int\delta_s r(t)=0$ because both densities have mass one.
Thus Lemma~\ref{multiplier} applies.

By joint continuity of $\partial_s f_\eta$, after shrinking $J$ if
necessary there is a modulus $\omega_f(\varepsilon)\to0$ such that
\[
 |\partial_s f_\eta(1+z)-a_\eta|
       \leq\omega_f(\varepsilon)\quad
       (\eta\in J,\ |z|\leq\varepsilon).
\]
If $\|r\|_\infty,\|g\|_\infty\leq\varepsilon$, the mean-value
formula and difference of squares give pointwise
\begin{align*}
 |\delta_s H|&\leq\omega_f(\varepsilon)|\delta_s r|
                               +\varepsilon|\delta_s g|,\\
 |\delta_s J|&\leq\varepsilon|\delta_s r|
                               +\varepsilon|\delta_s g|.
\end{align*}
In the second line use
$r(t+s)g(t+s)-r(t)g(t)
 =r(t+s)\delta_s g+g(t)\delta_s r$.
Minkowski's inequality and $\|z\|_2$ for each component bounded by
the combined two-component norm show
\[
 \bigl(\|\delta_s H\|_2^2+\|\delta_s J\|_2^2\bigr)^{1/2}
 \leq[\omega_f(\varepsilon)+3\varepsilon]
       \bigl(\|\delta_s g\|_2^2+\|\delta_s r\|_2^2\bigr)^{1/2}.
\]
Choose $0<\varepsilon_0<1$ so small that
$K[\omega_f(\varepsilon_0)+3\varepsilon_0]<1$.
Combining this inequality with \eqref{multbound} forces
$\delta_s g=\delta_s r=0$ in $L^2$ for every $s$.
The fields are continuous, so their differences vanish pointwise.
Hence $Dv$ and $m$ are independent of time. It follows that
$v(t,x)=V(x)+c(t)$. The HJB equation makes $c'(t)$ a constant,
since its other terms are independent of time. Periodicity makes that
constant zero, so $v$ itself is stationary.

For the final assertion, assume a genuine bifurcation sequence converges
in the stated rescaled $C^{1,2}$ topology with $P_n\to P_0>0$.
Its parameters eventually lie in $J\times I$, and its density and
gradient eventually satisfy the preceding smallness condition. Under
nonresonance it would therefore consist of stationary solutions, a
contradiction. This proves the claimed necessary resonance.
\end{proof}

The finite-period assumption matters: this estimate is uniform only for
periods in a neighborhood of $P_0>0$. It makes no claim about branches
with $P_n\to\infty$ or $P_n\to0$. A nontrivial spatial stationary
kernel cannot obstruct the proof because taking temporal differences
removes precisely the temporal zero modes.

\section{Remaining gap}
Theorem~\ref{exclude} is a nonlinear exclusion theorem valid for every
period, and \eqref{resonance} gives the full homogeneous linear list.
Theorem~\ref{nonlinearresonance} upgrades the earlier tangent-based test
to every genuine finite-period bifurcation sequence, including sequences
with a stationary or vanishing tangent. We have not proved nonlinear sufficiency
at every resonance, treated spatially dependent couplings or general
Hamiltonians, or continued the known periodic branches globally.
Those are substantial parts of the catalogue target, which remains
unresolved by this attempt.

\section{Completion Estimate}
35\%

This is a heuristic assessment of the full catalogue target. The version gives nonlinear exclusion regions and a necessary resonance classification for every finite-period bifurcation sequence in the homogeneous quadratic family. It does not prove nonlinear sufficiency at each resonance, classify spatially dependent couplings or general Hamiltonians, or continue periodic branches globally.

\begin{thebibliography}{9}
\bibitem{CN} M. Cirant and L. Nurbekyan,
\emph{The variational structure and time-periodic solutions for mean-field games systems},
Minimax Theory and its Applications 3 (2018), 227--260.
\url{https://arxiv.org/pdf/1804.08943v1}. Section 6, Theorem 6.1 and Remark 6.4.
\bibitem{CD} P. Cardaliaguet and F. Delarue,
\emph{Selected topics in mean field games}, Proceedings of the ICM 2022,
vol. 5, Section 3.3. DOI \href{https://doi.org/10.4171/ICM2022/17}{10.4171/ICM2022/17}.
\url{https://ems.press/content/book-chapter-files/33265}.
\end{thebibliography}
\end{document}
